\documentclass[11pt,a4paper]{article}
\usepackage[margin=25mm]{geometry}
\usepackage[T1]{fontenc}
\usepackage[utf8]{inputenc}
\usepackage{lmodern,microtype,amsmath,amssymb,amsthm,booktabs,longtable,enumitem}
\usepackage[round,sort]{natbib}
\usepackage{xurl}
\usepackage[colorlinks=true,allcolors=blue]{hyperref}
\setlength{\emergencystretch}{2em}
\numberwithin{equation}{section}
\newtheorem{theorem}{Theorem}[section]
\newtheorem{proposition}[theorem]{Proposition}
\newtheorem{lemma}[theorem]{Lemma}
\newtheorem{corollary}[theorem]{Corollary}
\newtheorem{conjecture}[theorem]{Conjecture}
\newtheorem{maintheorem}[theorem]{Theorem}
\theoremstyle{definition}
\newtheorem{definition}[theorem]{Definition}
\newtheorem{fact}{Standard fact}
\renewcommand{\thefact}{I\arabic{fact}}
\theoremstyle{remark}
\newtheorem{remark}[theorem]{Remark}
\newcommand{\Fp}{\mathbb F_p}
\newcommand{\Fq}{\mathbb F_{p^2}}
\newcommand{\Fb}{\overline{\mathbb F}_p}
\newcommand{\Z}{\mathbb Z}
\newcommand{\Q}{\mathbb Q}
\newcommand{\Res}{\operatorname{Res}}
\newcommand{\Aut}{\operatorname{Aut}}
\newcommand{\Stab}{\operatorname{Stab}}
\newcommand{\den}{\operatorname{den}_{\mathrm{red}}}
\newcommand{\leg}[2]{\left(\frac{#1}{#2}\right)}
\newcommand{\Lp}[1]{\left(\frac{#1}{p}\right)}
% Set this to the public evidence repository URL or DOI landing page before submission.
\newcommand{\DataRepositoryURL}{https://github.com/bbrhuft/fricke-signed-residues/tree/main/supplementary}

% ---------------------------------------------------------------------
% v2.0 (9 Nov 2026). Paper B of two. The combined manuscript v1.14 was split:
% this paper contains the Fricke application (old Sections 1.3-1.4, 4-8).
% The factorial evaluation and the signed discriminant-ratio identity are in
% the companion paper "A factorial supersingular congruence and signed
% discriminant ratios" and enter here only through Theorem 1.1.
% Statements and proofs are those of v1.14, except:
%  * Theorem 1.2(d) and Proposition 4.3 now state r_i W_i = beta(E) D_N(x)^k;
%    the further expression (N/p) c(E) c(E') was dropped, since it needs
%    Robert's isogeny identity and the constant c(E) of the companion paper;
%  * part (e) is deduced from Baker's congruence and the companion theorem.
% Open items for the author: search for "TODO(author)".
% ---------------------------------------------------------------------
\title{Supersingular poles and signed residues\\of prime-level Fricke Hauptmoduln}
\author{David Jordan}
\date{9 November 2026}
\begin{document}
\maketitle
\begin{abstract}
Let $N$ be one of the fifteen primes for which $X_0(N)^+$ has genus zero,
let $T_N=q^{-1}+O(q)$ be its Hauptmodul, and let $p\ge5$ be a prime with
$p\ne N$. We determine the poles and signed residues of the rational
function representing $T_N|U_p$ modulo $p$. Its poles are precisely the
supersingular values with effective ramification index one, and each
residue multiplied by the weighted product of differences from all
other supersingular values equals $-\left(\frac{-N}{p}\right)$.
Consequently the supersingular values and their indices determine the
reduced function. The sign rests on Baker's supersingular congruence and
on the signed discriminant-ratio identity of a companion paper.
The indices also describe node thicknesses in the Deligne--Rapoport
Fricke quotient. Exact computations provide independent finite
certificates for the residue formula at all fifteen levels.
\end{abstract}
\noindent\textbf{Keywords:} supersingular elliptic curves; Fricke groups; Hauptmoduln; Hecke operators; eta quotients; congruences.

\section{Introduction}\label{sec:intro}
This paper concerns the reduction of $U_p$ applied to a
Fricke Hauptmodul. At each of the fifteen prime levels with genus-zero
Fricke quotient, we determine the pole divisor and signed residues of
the reduced function. The ramification indices of the quotient remove
certain supersingular points from the pole divisor, but those points
remain in the weighted residue product. The formula reconstructs the
entire reduced function from the supersingular values and their indices.

The sign in the residue formula comes from two supersingular identities,
recalled in Theorem~\ref{thm:input}: Baker's congruence \citep{Baker98}
and the signed discriminant-ratio identity proved in the companion paper
\citep{JordanFactorial}.

\subsection{Notation}
Let $p\ge5$ be prime and put
\[
h=(p-1)/2,\qquad k=(p^2-1)/24,\qquad
A=\overline E_{p-1},\qquad B=\overline E_{p+1},\qquad
\Delta=\eta^{24}.
\]
Here $E_w$ is the normalized Eisenstein series of weight $w$, a bar
denotes reduction modulo $p$, and $A$ is the Hasse invariant. The
$q$-expansions of $A$ and $B$ are $1$ and $E_2$ modulo $p$,
respectively. We write $q=e^{2\pi i\tau}$, $\theta=q\,d/dq$,
$\Fp=\mathbb F_p$, $\Fq=\mathbb F_{p^2}$ and
$\Fb=\overline{\mathbb F}_p$.
For $f=\sum c_nq^n$, define $f|U_p=\sum c_{pn}q^n$.
The symbol $\left(\frac{a}{p}\right)$ denotes the Legendre symbol.
When a modular form is evaluated at an elliptic curve, a nonzero
invariant differential is understood. Ratios of equal weight are
independent of its choice. In particular, set
\[
\beta(E):=\frac{B(E,\omega)^{p-1}}{\Delta(E,\omega)^{2k}}.
\]

\subsection{Two supersingular identities}
For a prime $N\ne p$, put $D_N=\Delta(\tau)/\Delta(N\tau)$. It is a ratio
of two forms of weight $12$ on $\Gamma_0(N)$ with integral $q$-expansions,
neither of which vanishes on the non-cuspidal curve $Y_0(N)$ over
$\Z[1/N]$. It is therefore a unit on $Y_0(N)$, and we use the same
notation for its reduction modulo $p$. A description of $D_N$ in terms of
the isogeny $E\to E/C$ is given in \citet{JordanFactorial}.

% TODO(author): the theorem number for Baker is that of the author's posted
% version of the paper; check it against the journal version.
\begin{theorem}[supersingular inputs]\label{thm:input}
Let $p\ge5$ and $N\ne p$ be primes.
\begin{enumerate}[label=\textup{(\roman*)}]
\item \textup{\citep[Theorem~1.4]{Baker98}.} For every supersingular elliptic
curve $E/\Fb$,
\[
\beta(E)=-\Lp{-1}.
\]
\item \textup{\citep[Theorem~1.3]{JordanFactorial}.} At every supersingular
point $x$ of $Y_0(N)_{\Fb}$,
\[
D_N(x)^{(p^2-1)/24}=\Lp N.
\]
\end{enumerate}
\end{theorem}

Part (ii) refines the identity $D_N(x)^{(p^2-1)/12}=1$ of
\citet[Theorem~1.2]{Ono25}. It is deduced in \citet{JordanFactorial} from
the evaluation $B(E,\omega)^h=\Lp6\,h!\,\Delta(E,\omega)^k$ on the
supersingular locus, which also recovers part (i).

\subsection{The residue formula}
Throughout,
\[
N\in\mathcal N:=\{2,3,5,7,11,13,17,19,23,29,31,41,47,59,71\},
\]
and $p\ge5$ is prime with $p\ne N$.
These are precisely the prime levels for which $X_0(N)^+$ has genus
zero \citep[\S3, Corollary]{Ogg75}. Let $T_N$ denote its unique
Hauptmodul with expansion $q^{-1}+O(q)$ at the single cusp.
Lemma~\ref{lem:coordinate} constructs this coordinate over
$\Lambda_N:=\Z[1/(6N)]$ and proves that it remains a coordinate
at every admissible prime. The standard moonshine identification
relates these functions to prime-order Fricke McKay--Thompson series
\citep{CN79,CMT19}; the proof below uses the modular curves.

For the explicit applications at
$N\in\mathcal N_0:=\{2,3,5,7,13\}$ we also use
\[
a_N=\frac{24}{N-1},\qquad b_N=N^{12/(N-1)},\qquad
u_N=\left(\frac{\eta(\tau)}{\eta(N\tau)}\right)^{a_N},\qquad
T_N=u_N+a_N+b_N/u_N.
\]
The pairs $(a_N,b_N)$ are $(24,4096),(12,729),(6,125),(4,49),(2,13)$.
Here $u_N$ is a coordinate on $X_0(N)$ and $w_N$ sends it to $b_N/u_N$.
This eta presentation is not used to define $T_N$ at the ten other levels.

Write $T=T_N$ and
$\Delta_N=\Delta(\tau)\Delta(N\tau)$. The level $N$ and the
reduction prime $p$ play distinct roles: $p$ need not belong to
$\mathcal N$, and $p=N$ is excluded.
At the five eta-coordinate levels, $D_N=u_N^{N-1}$, so
Theorem~\ref{thm:input}(ii) gives
\[
u_N(x)^{(p^2-1)/a_N}=\Lp N
\]
at every supersingular point $x$ of $Y_0(N)$.

The following weights distinguish the supersingular points that
contribute poles from those at which ramification cancels a pole.

Let $t_1,\dots,t_s\in\Fb$ be the distinct supersingular values of $T$. For a point
$x=(E,C)$ of $X_0(N)_{\Fb}$ over $t_i$ put
\[
m_x=\bigl|\Aut(E,C)/\{\pm1\}\bigr|,\qquad e_i=m_x\,\bigl|\Stab_{\langle w_N\rangle}(x)\bigr|.
\]
Call $t_i$ \emph{thin} if $e_i=1$ and \emph{thick} if $e_i>1$; the terms refer to the
node-thickness interpretation of Section~\ref{sec:geometry}. Put
$D_{\rm thin}(T)=\prod_{e_i=1}(T-t_i)$. With $R_{N,p}$ the rational function of
Theorem~\ref{thm:main}(a) below, let, for thin $t_i$,
\[
r_i=\Res_{T=t_i}R_{N,p}(T),\qquad W_i=\prod_{j\ne i}(t_i-t_j)^{(p+1)/e_j}.
\]
Section~\ref{sec:weights} shows that $e_i$ does not depend on the representative of the Fricke
orbit, that it is the effective (tame) inertia of the quotient map at $x$ (not the order of a group
containing $\pm1$), and that every exponent $(p+1)/e_j$ is an integer.

Only thin values give poles of $R_{N,p}$, and every such pole is simple
(Theorem~\ref{thm:main}(b)).

\begin{maintheorem}[Fricke residue formula]\label{thm:main}
Let $N\in\mathcal N$ and let $p\ge5$ be a prime with $p\ne N$.
\begin{enumerate}[label=\textup{(\alph*)}]
\item \textbf{(Rationality.)} There is a unique $R_{N,p}\in\Fp(T)$ with $\overline{T|U_p}=R_{N,p}(\overline T)$. It vanishes at $T=\infty$.
\item \textbf{(Pole support.)} $R_{N,p}=P/D_{\rm thin}$ with $P\in\Fp[T]$,
$\deg P<\deg D_{\rm thin}$ and $\gcd(P,D_{\rm thin})=1$. Every pole is simple and has a nonzero residue. In particular $\den R_{N,p}=D_{\rm thin}$.
\item \textbf{(Residue as a Hasse derivative.)} If $t_i$ is thin and $x$ lies over $t_i$, then
\[
r_i=\frac{144\,(\theta T)(x)^{p+1}}{B(x)^2},\qquad \theta=q\frac{d}{dq},
\]
a ratio of two mod-$p$ modular forms of weight $2(p+1)$, evaluated at $x$.
\item \textbf{(Product.)} For thin $t_i$ and $x=(E,C)$ over $t_i$,
\[
r_iW_i=\frac{B^{p-1}}{\Delta_N^{\,k}}(x)=\beta(E)\,D_N(x)^k.
\]
\item \textbf{(Signed weighted residue formula.)} For every thin $t_i$,
\[
\boxed{\;r_iW_i=-\Lp{-N}\;}\qquad\text{equivalently}\qquad
r_i=-\Lp{-N}\prod_{j\ne i}(t_i-t_j)^{-(p+1)/e_j}.
\]
In particular $r_i\in\Fp^\times$ and $W_i\in\Fp^\times$, also when $t_i\notin\Fp$.
\item \textbf{(Hecke form.)} With $\Psi_{N,p}(X,Y)=(X^p-Y)(X-Y^p)+pH_{N,p}(X,Y)$ as in
Proposition~\ref{prop:Psi}, $\overline H_{N,p}(t^p,t)=\Lp{-N}W_i^{-1}$ if $t=t_i$ is thin, and
$\overline H_{N,p}(t^p,t)=0$ for every other $t\in\Fq$.
\end{enumerate}
\end{maintheorem}

\begin{corollary}[reconstruction]\label{cor:reconstruction}
The supersingular values and their indices determine the entire reduced function:
\begin{equation}\label{eq:partial-fractions}
R_{N,p}(T)=-\Lp{-N}\sum_{e_i=1}\frac{1}{W_i(T-t_i)}.
\end{equation}
Equivalently, if $D_{\rm thin}=\prod_f f$ is its factorization into monic
irreducibles over $\Fp$, then
\begin{equation}\label{eq:factor-reconstruction}
R_{N,p}(T)=\sum_{f\mid D_{\rm thin}} c_f\frac{f'(T)}{f(T)},\qquad
c_f=-\Lp{-N}W_i^{-1}\in\Fp^\times\quad\text{for any root }t_i\text{ of }f.
\end{equation}
Both sums are empty when $D_{\rm thin}=1$.
\end{corollary}
\begin{proof}
Subtracting the principal parts leaves a polynomial, which is zero because the
function and each summand vanish at infinity. Frobenius-conjugate poles have
conjugate residues; the residues lie in $\Fp$ by Theorem~\ref{thm:main}(e), so
they are equal on each irreducible factor. The identity
$f'/f=\sum_{f(t_i)=0}(T-t_i)^{-1}$ gives the second formula.
\end{proof}


The residues determine the reduced function because it has only simple
poles and vanishes at the cusp. Corollary~\ref{cor:reconstruction}
therefore reconstructs $R_{N,p}$ from the supersingular values and their
indices alone. Remark~\ref{rem:datasize} compares the size of these data
with that of a $q$-expansion certificate.
The same effective indices are interpreted as node thicknesses in
Section~\ref{sec:geometry}.

\subsection{Consequences of the Fricke application}\label{subsec:consequences}
The reconstruction in Corollary~\ref{cor:reconstruction} gives concrete formulas
once the supersingular values and their effective indices are known.
At the five levels $N\in\mathcal N_0$, the connection to Nakaya's
squarefree Fricke supersingular polynomials is a change of coordinate:
with the normalizations in \citet[\S4]{Nakaya19},
\begin{equation}\label{eq:nakaya-coordinate}
\mathcal S_{N,p}(T):=\prod_i(T-t_i)=ss_p^{(N*)}(T+c_N),\qquad
\begin{array}{c|rrrrr}
N&2&3&5&7&13\\\hline
c_N&104&42&16&9&0
\end{array}
\end{equation}
For $N=2,3,5,7$, Nakaya's printed eta formulas have additive
constants $128,54,22,13$, giving the shifts displayed above. At level
$13$ his convention is zero constant term in the $q$-expansion; this
identifies the normalized Hauptmodul with $T_{13}$ and gives $c_{13}=0$.
All constants are reduced modulo $p$; $\mathcal S_{N,p}$ is distinct from the
weighted polynomial $S$ of Proposition~\ref{prop:S}. Thus
$D_{\rm thin}=\mathcal S_{N,p}/\prod_{e_i>1}(T-t_i)$, while the omitted
values still occur in the weights $W_i$.
For these five levels the exceptional values can also be specified
by a polynomial independent of $p$.
Let $\mathfrak h_N$ be the elliptic-value polynomial in Table~\ref{tab:elliptic-values}.
Lemma~\ref{lem:elliptic-specialization} and Theorem~\ref{thm:main}(b) give
\begin{equation}\label{eq:elliptic-gcd}
D_{\rm thin}(T)=\frac{\mathcal S_{N,p}(T)}
{\gcd\bigl(\mathcal S_{N,p}(T),\overline{\mathfrak h}_N(T)\bigr)},
\end{equation}
where the gcd is monic. Thus the denominator is obtained from Nakaya's
supersingular polynomial by removing its common factors with the reduction
of a fixed characteristic-zero elliptic-value polynomial.

At level two, put
$\varepsilon_p=(1-\left(\frac{-1}{p}\right))/2$ and
$\nu_p=(1-\left(\frac{-2}{p}\right))/2$.
The level-two indices described in Section~\ref{sec:geometry} give
\begin{equation}\label{eq:leveltwo-denominator}
D_{2,p}(T)=\frac{ss_p^{(2*)}(T+104)}
{(T+104)^{\varepsilon_p}(T-152)^{\nu_p}}.
\end{equation}
The two removed values correspond to $0$ and $256$ in Nakaya's coordinate;
they are distinct for $p\ge5$.

\begin{corollary}[level-two degree and vanishing]\label{cor:leveltwo-vanishing}
For every prime $p\ge5$,
\begin{equation}\label{eq:leveltwo-degree}
\deg D_{2,p}=\frac{p-4+\left(\frac{-1}{p}\right)+2\left(\frac{-2}{p}\right)}8.
\end{equation}
In particular $T_{2A}|U_p\equiv0\pmod p$ if and only if $p=5$ or $7$.
\end{corollary}
\begin{proof}
The two possible thick values have indices four and two, and supersingularity
indicators $\varepsilon_p$ and $\nu_p$, respectively. Corollary~\ref{cor:degree}
therefore gives
$\deg D_{2,p}=(p-1)/8-\varepsilon_p/4-\nu_p/2$, which is
\eqref{eq:leveltwo-degree}. This degree is zero at $p=5,7$ and is at least
$(p-7)/8>0$ for $p\ge11$. Theorem~\ref{thm:main}(b) proves the last assertion.
\end{proof}

\citet[Lemma~4.8 and Table~4.3]{CMT19} establish $p$-adic
annihilation of $T_{2A}$ at $p=5,7$. The corollary specifies the exact
first-step modulo-$p$ vanishing criterion. Table~\ref{tab:leveltwo-examples}
illustrates the nonzero cases; at $p=17$, for example,
$R_{2,17}=D_{2,17}'/D_{2,17}$ and both residues are one.

\begin{table}[htbp]\centering
\begin{tabular}{rll}\toprule
$p$ & $D_{2,p}(T)$ & $P_{2,p}(T)$\\\midrule
11 & $T-5$ & $1$\\
13 & $T-3$ & $-2$\\
17 & $T^2-T-8$ & $2T-1$\\
29 & $T^3-12T^2+12T-6$ & $4T^2-5T-8$\\\bottomrule
\end{tabular}
\caption{Level-two examples, with $R_{2,p}=P_{2,p}/D_{2,p}$.
Coefficients use the unique integer representatives in $(-p/2,p/2)$.
These entries are included in the certificates of Section~\ref{sec:comp}.}
\label{tab:leveltwo-examples}
\end{table}

\subsection{Relation to earlier work}
The sign input, Theorem~\ref{thm:input}, combines \citet{Baker98} with
\citet{JordanFactorial}; the remaining ingredients are as follows.

The derivative argument in Section~\ref{sec:kz} follows the level-one
method presented by \citet[\S10]{KZ98}, who trace it to Koike's work
on Atkin's congruences \citep{Koike73}. The period formula of
\citet{dS94} and Baker's congruence supply the level-one precedents
for the residue product and its sign; \citet{BL18} treat the
corresponding Legendre modular polynomials. The Fricke application
combines this derivative method with Theorem~\ref{thm:input} and
the effective ramification indices of the quotient.

Fricke supersingular polynomials are studied by \citet{Sakai11,Sakai14},
\citet{Nakaya19} and \citet{Morton23}. These papers treat
orthogonal-polynomial representations, factorization and resultant
identities. Formula~\eqref{eq:nakaya-coordinate} identifies the supersingular
values used here with that literature; the present residue formula additionally
assigns a nonzero coefficient to each thin pole and retains the ramified values
in its weighted product.

\citet[\S4.2, Table~1 and Appendix~C]{JST19} determine the elliptic values
of normalized Fricke Hauptmoduln and their associated class fields. Their
$j_N$ equals our $T_N$. Table~\ref{tab:elliptic-values} records their
polynomials for all fifteen prime levels. Their Schwarzian equation
detects the same local ramification indices that enter our Hasse-derivative
calculation; \eqref{eq:elliptic-gcd} expresses its elliptic-value locus as
the locus of possible supersingular pole cancellations. The signed residue
formula concerns the remaining thin values.

\citet[Proposition~3.1 and \S4]{DS26} use Deligne's supersingular
expansion of $pJ|U_p$ and prime-power level modular functions to
characterize the monstrous exponents. Here we determine the signed
residues on the fifteen prime-level genus-zero Fricke curves.

\citet[Theorem~1.1]{BGKMR26} study iterated $U_p$ congruences for
positive-weight meromorphic cusp forms on $\mathrm{SL}_2(\Z)$ whose
interior pole reduces to a supersingular point. Their $p$-adic
approximation by holomorphic cusp forms concerns a different setting
from our first-step modulo-$p$ pole divisor and signed residue formula.

\paragraph{Organization.}
Section~\ref{sec:imports} constructs the integral coordinate and proves
the residue criterion. Section~\ref{sec:weights} treats the effective
indices and the weighted supersingular polynomial. Section~\ref{sec:kz}
proves Theorem~\ref{thm:main}. Section~\ref{sec:geometry} interprets the
weights geometrically, and Section~\ref{sec:comp} records finite
certificates and exact-arithmetic evidence.

\paragraph{Acknowledgement of AI assistance.}
% TODO(author): confirm this statement (including the list of models) and
% check the target journal's AI policy.
OpenAI GPT-6.1 Sol and Claude Opus 5.5 assisted in developing arguments, checking
calculations and preparing this manuscript. The source computations
and proof checks are documented in the supplementary packages.
The author takes responsibility for the entire content of the paper.

\section{Standard inputs and the residue criterion}\label{sec:imports}

\begin{lemma}[integral Fricke coordinate]\label{lem:coordinate}
Let $N$ be prime with $g(X_0(N)^+)=0$, and put
$\Lambda_N=\Z[1/(6N)]$. The coarse quotient
\[
\mathcal X_N^+:=X_0(N)_{\Lambda_N}/\langle w_N\rangle
\]
is a smooth proper curve over $\Lambda_N$ with a single cusp section $c$.
There is a unique coordinate $T_N$ over $\Lambda_N$ whose expansion at
that section is $q^{-1}+O(q)$. It identifies $\mathcal X_N^+$ with
$\mathbb P^1_{\Lambda_N}$ and $c$ with infinity. In particular,
$T_N\in q^{-1}+q\Lambda_N[[q]]$ and its reduction is a coordinate
at every prime $p\ge5$ with $p\ne N$.
\end{lemma}
\begin{proof}
Away from $6N$ the coarse modular curve $X_0(N)$ is smooth and proper,
with geometrically connected fibres; its two disjoint rational cusp
sections are exchanged by $w_N$ \citep{DR73,KM85}.
The quotient exists and is proper over $\Lambda_N$. To see that it is flat,
work on an invariant affine open with coordinate ring $A$. Since $2$ is
invertible in $\Lambda_N$, the averaging map
\[
a\longmapsto \frac{a+w_N(a)}{2}
\]
makes the invariant ring $A^{\langle w_N\rangle}$ a direct summand of $A$.
As $A$ is flat over $\Lambda_N$, its direct summand is also flat. Since $2$ is invertible,
formation of the quotient commutes with base change. Each geometric
fibre is the quotient of a smooth curve by an involution, hence is
smooth; thus $\mathcal X_N^+$ is smooth over $\Lambda_N$.
Its fibres are geometrically connected and have constant genus,
which is zero by the assumption on the generic fibre.
The two cusps give one section $c$. The quotient map is an isomorphism
on the completed neighbourhood of either cusp onto that of $c$, so
$\widehat{\mathcal O}_{\mathcal X_N^+,c}=\Lambda_N[[q]]$.

Write $f:\mathcal X_N^+\to\operatorname{Spec}\Lambda_N$.
On every fibre, $\mathcal O(c)$ has degree one and $H^1(\mathcal O)=0$.
Cohomology and base change therefore make $f_*\mathcal O(c)$ locally
free of rank two, and the principal-part sequence is exact:
\[
0\longrightarrow\Lambda_N\longrightarrow
H^0(\mathcal X_N^+,\mathcal O(c))
\longrightarrow\Lambda_N\,q^{-1}\longrightarrow0.
\]
The last identification uses the parameter $q$ to trivialize the
normal line of the section. Choose a section with principal part
$q^{-1}$ and subtract its constant coefficient. This gives $T_N$;
two such choices differ by a global constant, proving uniqueness.
The sections $1,T_N$ form a basis on every fibre and generate
$\mathcal O(c)$ there, so they give an isomorphism to
$\mathbb P^1_{\Lambda_N}$, with the asserted cusp and coordinate.
\end{proof}

The lemma uses genus zero only for the Fricke quotient. Neither its
proof nor the arguments below require $X_0(N)$ itself to have genus zero.


% TODO(author): the references below are given at section or chapter level;
% add theorem-level pinpoints after checking them against the sources.
We use the following standard results for prime-to-$p$ level structures.
The $q$-expansion principle is \citet[\S1.6]{Katz73}. The Hasse invariant,
its $q$-expansion and Deligne's congruence $A\equiv E_{p-1}\pmod p$ are
treated in \citet[\S\S2.0--2.1]{Katz73}. That the Hasse invariant has simple
zeros, precisely at the supersingular points of a fine moduli cover, is
Igusa's theorem \citep[Ch.~12]{KM85}. The Kodaira--Spencer isomorphism,
under which a weight-two form $f$ corresponds to $f\,dq/q$, is in
\citet[Appendix~1]{Katz73}; see also \citet{Gross90}. In particular,
$\theta T$ represents $dT$. We record the statements used repeatedly.
\begin{fact}\label{fact:modp}\leavevmode
\begin{enumerate}[label=\textup{(\roman*)}]
\item The $q$-expansion principle holds at prime-to-$p$ level. The forms
$A$ and $B$ have $q$-expansions $1$ and $E_2$ modulo $p$, respectively.
\item The Hasse invariant $A$ has simple zeros precisely at the
supersingular points on a fine moduli cover.
\item For $p\ge5$ the stabilizers relevant to $X_0(N)$ and its Fricke
quotient are tame, and these coarse quotients commute with reduction.
The supersingular values of $T$ lie in $\Fq$.
\item Kodaira--Spencer identifies a weight-two form $f$ with
$f\,dq/q$; in particular $\theta T$ represents $dT$.
\end{enumerate}
\end{fact}
The field-of-definition assertion in (iii) follows by choosing a model
with $p^2$-Frobenius $[-p]$ as follows; every cyclic subgroup of order
$N$ is then rational as a subgroup scheme.

\begin{lemma}[Deuring models]\label{lem:deuring}
Let $E_1/\Fb$ be supersingular, $p\ge5$. Then $E_1$ has a model over $\Fq$ whose $p^2$-power
Frobenius endomorphism is $[-p]$. For such a model every finite subgroup of order prime to $p$ is
$\Fq$-rational, and $E_1/K$ again has Frobenius $[-p]$.
\end{lemma}
\begin{proof}
A supersingular curve admits a model over $\Fq$
\citep{Deuring41,Waterhouse69}. For any such model $E/\Fq$, the $p$-power Frobenius $\mathrm{Fr}:E\to E^{(p)}$ and the
Verschiebung $V:E^{(p)}\to E$ are both purely inseparable of degree $p$ (the Hasse invariant
vanishes), so $\mathrm{Fr}^{(p)}:E^{(p)}\to E^{(p^2)}=E$ and $V$ differ by an automorphism $u$ of $E$:
$\mathrm{Fr}^{(p)}=u\circ V$, whence
$\pi_{p^2}=\mathrm{Fr}^{(p)}\circ\mathrm{Fr}=u\circ V\circ\mathrm{Fr}=u\circ[p]$. Thus
$\pi_{p^2}=p\,u$ with $u\in\Aut(E)$. If $j\ne0,1728$ then $u=\pm1$ and the quadratic twist changes
the sign; if $j=0$ (resp.\ $1728$) the curve $y^2=x^3+1$ (resp.\ $y^2=x^3+x$) over $\Fp$ has
trace $0$, so $\pi_{p^2}=\pi_p^2=[-p]$ \citep{Deuring41,Waterhouse69}; see also
\citet[Lemma~4.3]{Ono25}. Since $[-p]$ is central, Galois preserves every subgroup, and the quotient
isogeny $\phi$ satisfies $\pi'\phi=\phi\pi=\phi[-p]=[-p]\phi$, so $\pi'=[-p]$.
\end{proof}

\begin{proposition}\label{prop:Psi}
Let $\Delta_p$ be the set of integral matrices of determinant $p$ with lower-left entry divisible by
$N$. There is a unique $\Psi_{N,p}\in\Lambda_N[X,Y]$ with
\[
\Psi_{N,p}(X,T(\tau))=(X-T(p\tau))\prod_{b=0}^{p-1}\Bigl(X-T\bigl(\tfrac{\tau+b}{p}\bigr)\Bigr).
\]
It is monic of degree $p+1$ in each variable, symmetric, and satisfies
$\Psi_{N,p}\equiv(X^p-Y)(X-Y^p)\pmod p$.
\end{proposition}
\begin{proof}
The matrices $\mathrm{diag}(p,1)$ and $\begin{pmatrix}1&b\\0&p\end{pmatrix}$, $0\le b<p$, represent
$\Gamma_0(N)\backslash\Delta_p$. Conjugation by $W_N=\begin{pmatrix}0&-1\\N&0\end{pmatrix}$ maps
$\begin{pmatrix}a&b\\c&d\end{pmatrix}$ to $\begin{pmatrix}d&-c/N\\-Nb&a\end{pmatrix}$, so it preserves
$\Delta_p$. Hence the elementary symmetric functions of the roots are $\Gamma_0(N)^+$-invariant,
holomorphic on the upper half-plane, with a pole only at the single cusp.
Their coefficients initially lie in $\Lambda_N[\zeta_p]((q^{1/p}))$.
Translation of $\tau$ removes fractional powers, and cyclotomic Galois automorphisms
permute the roots, so the expansions lie in $\Lambda_N((q))$.
Subtracting powers of $T=q^{-1}+O(q)$ then puts them in $\Lambda_N[T]$. Modulo $1-\zeta_p$ all $T((\tau+b)/p)$ coincide with $T(q^{1/p})$, whose $p$-th power is $T(q)$, and
$T(p\tau)\equiv T^p$; this gives the congruence. For irreducibility, the reduction of $\Gamma_0(N)$ modulo $p$ is
$\mathrm{SL}_2(\Fp)$, by the Chinese remainder theorem and the surjectivity of reduction
for $\mathrm{SL}_2(\Z)$. It acts transitively on the $p+1$ cyclic subgroups of order $p$.
The $p+1$ roots displayed above are distinct: their leading terms are $q^{-p}$ and
$\zeta_p^{-b}q^{-1/p}$ for $0\le b<p$. They therefore form one orbit, so
$\Psi_{N,p}(X,T)$ is irreducible over $\mathbb C(T)$.
The pole order of the $j$th elementary symmetric function is at most
$p+(j-1)/p\le p+1$; the constant coefficient has pole order $p+1$.
Thus the polynomial has degree $p+1$ in $Y$ as well as in $X$.
For symmetry, $\mathrm{diag}(1,p)\in\Delta_p$ gives
$\Psi_{N,p}(T(\tau),T(p\tau))=0$ by evaluating the defining relation at $p\tau$.
Irreducibility identifies the reversed polynomial with the original
up to a rational scalar. The constant coefficient, as a polynomial in
$X$, has leading term $q^{-(p+1)}$ with coefficient one, so
$\Psi_{N,p}$ is monic of degree $p+1$ in $Y$. Comparing the
$X^{p+1}$ coefficients in the two polynomials makes that scalar one.
Thus $\Psi_{N,p}(Y,X)=\Psi_{N,p}(X,Y)$.
\end{proof}

Define $H=H_{N,p}\in\Lambda_N[X,Y]$ by $\Psi_{N,p}=(X^p-Y)(X-Y^p)+pH$; $H$ is symmetric. Put
\[
\psi=\frac{T(q)^p-T(q^p)}{p}\in\Lambda_N((q)),
\]
which lies in this ring because
$T(q^p)\equiv T(q)^p\pmod{p\Lambda_N}$.

\begin{proposition}[residue criterion]\label{prop:res}
In $\Fp((q))$,
\[
\overline\psi=\frac{\overline H(T^p,T)}{T^{p^2}-T},\qquad
\overline{T|U_p}=\overline Q(T)+\overline\psi,
\]
where $\overline Q\in\Fp[T]$ is defined by $-[X^p]\Psi_{N,p}=Y^p+pQ(Y)$. Hence
$R_{N,p}=\overline Q(T)+\overline H(T^p,T)/(T^{p^2}-T)$ exists and is unique. It vanishes at
$\infty$, its poles are simple and lie in $\Fq$, and for $t\in\Fq$
\[
\Res_{T=t}R_{N,p}=-\overline H(t^p,t).
\]
\end{proposition}
\begin{proof}
Substitute $X=T(q^p)=T^p-p\psi$ and $Y=T$ into
$\Psi_{N,p}=0$. Since $X-Y^p=-p\psi$, division by $p$ and reduction give
$-\overline\psi\,(T^{p^2}-T)+\overline H(T^p,T)=0$. The sum of the roots of $\Psi_{N,p}(X,T)$ is
$T(q^p)+p\,(T|U_p)=T^p+pQ(T)$, so $T|U_p=Q(T)+\psi$. The polynomial $T^{p^2}-T$ is squarefree, its
roots are exactly $\Fq$, and its derivative is $-1$. $R_{N,p}$ vanishes at $\infty$ because
$T|U_p\in q\Lambda_N[[q]]$ (the constant term of $T$ is $0$). Uniqueness holds because $\overline T$ is
transcendental over $\Fp$.
\end{proof}

\begin{remark}[Faber polynomials and replication]\label{rem:faber}
The polynomial
\[
\mathcal F_{N,p}(X):=-[Z^p]\Psi_{N,p}(Z,X)=X^p+pQ(X)
\]
is the degree-$p$ Faber polynomial of $T$: it is the unique monic
polynomial for which $\mathcal F_{N,p}(T(q))=q^{-p}+O(q)$. Indeed, the
sum-of-roots identity used above gives
\[
\mathcal F_{N,p}(T(q))=T(q^p)+p(T|U_p)(q).
\]
Under the standard McKay--Thompson identification of $T$, this is the
prime-to-$N$ replication identity; see \citet[\S8]{CN79} for the
replication terminology. The polynomial $Q$ is the correction
$(\mathcal F_{N,p}-X^p)/p$, rather than the Faber polynomial itself.
Proposition~\ref{prop:res} describes its reduction together with the
first Kronecker correction.
\end{remark}

\begin{corollary}\label{cor:rFp}
For every $t\in\Fq$, $\Res_{t}R_{N,p}\in\Fp$ and $\Res_{t^p}R_{N,p}=\Res_tR_{N,p}$.
\end{corollary}
\begin{proof}
$(\Res_tR)^p=\Res_{t^p}R=-\overline H(t^{p^2},t^p)=-\overline H(t,t^p)=-\overline H(t^p,t)=\Res_tR$, using
$t^{p^2}=t$, the $\Fp$-rationality of $R$ and the symmetry of $H$.
\end{proof}

Proposition~\ref{prop:res} gives Theorem~\ref{thm:main}(a) and the a~priori bound
$\den R_{N,p}\mid T^{p^2}-T$ used by the certificates in Section~\ref{sec:comp}.

\section{Weights and the weighted supersingular polynomial}\label{sec:weights}

\begin{lemma}[effective inertia]\label{lem:inertia}
Let $x=(E,C)$ be a non-cuspidal point of $X_0(N)_{\Fb}$, supersingular or ordinary. Let $G_x$ be the group of automorphisms
of $(E,C)$ together with the isomorphisms $(E,C)\to w_N(E,C)=(E/C,E[N]/C)$ when $w_N(x)=x$. It is a group once
$w_N^2(E,C)\cong(E,C)$ is identified via $[N]$.
\begin{enumerate}[label=\textup{(\roman*)}]
\item $|G_x/\{\pm1\}|=m_x|\Stab_{\langle w_N\rangle}(x)|=e_x$. The value is the same at $x$ and at $w_N(x)$.
\item $\{\pm1\}$ is exactly the kernel of the action of $G_x$ on the deformation space of $x$. Since
$e_x\in\{1,2,3,4,6\}$ is prime to $p$, $G_x/\{\pm1\}$ is cyclic and acts faithfully on the tangent line.
Hence, for a local parameter $z$ of a fine moduli cover at a point over $x$, $T-t=z^{e_x}\cdot(\text{unit})$.
\item If $x$ is supersingular, $(p+1)/e_x\in\Z$.
\end{enumerate}
\end{lemma}
\begin{proof}
(i) If $w_N(x)=x$, the isomorphisms $(E,C)\to w_N(E,C)$ form a torsor under $\Aut(E,C)$, so
$|G_x|=2|\Aut(E,C)|$. Conjugation by the isogeny $E\to E/C$ identifies $\Aut(E,C)$ with
$\Aut(E/C,E[N]/C)$, so $m_x=m_{w_N(x)}$.
(ii) $[-1]$ acts trivially on deformations. A Fricke-type element acts
on the deformation functor by
$(\widetilde E,\widetilde C)\mapsto
(\widetilde E/\widetilde C,\widetilde E[N]/\widetilde C)$,
followed by the chosen identification of $w_N(E,C)$ with $(E,C)$.
If its action were trivial, the universal deformation would admit an
isomorphism $\widetilde E/\widetilde C\cong\widetilde E$.
Composing with the quotient map would give a degree-$N$ endomorphism
of its generic fibre, whose endomorphism ring is $\Z$; this is impossible
since $N$ is not a square. Similarly, a nontrivial automorphism other
than $\pm1$ acting trivially would extend to the universal deformation
and contradict that generic endomorphism ring. The same argument works at ordinary points. A finite group of order prime to $p$ acting faithfully on $k[[z]]$ is cyclic and
acts faithfully on $(z)/(z^2)$, so the invariant ring is $k[[z^{e}]]$.
(iii) $e=3$ or $6$ forces $m_x=3$, so $j=0$ is supersingular and $p\equiv2\pmod3$; $e=4$ forces
$m_x=2$, so $j=1728$ is supersingular and $p\equiv3\pmod4$. For $p\ge5$, $m_x\in\{1,2,3\}$.
\end{proof}

For explicit calculations at $N\in\mathcal N_0$, write $u=u_N$.
Then $m_x$ can be read off from the multiplicity $\mu$ of $u(x)$ as a root of
$P_N(u)-j\,u^N$: $m_x=3/\mu$ over $j=0$, $m_x=2/\mu$ over $j=1728$, and $m_x=1$ (with $\mu=1$) otherwise.
Here $\mu$ is the ramification index of the coarse map $X_0(N)\to X(1)$.
The Fricke stabilizer has order two exactly when $u(x)^2=b_N$.

\begin{lemma}[elliptic values and specialization]\label{lem:elliptic-specialization}
Let $N\in\mathcal N$ and let
$\mathfrak h_N(T)=\prod_{\tau}(T-T_N(\tau))$, where $\tau$ runs over the
inequivalent elliptic points of $\Gamma_0(N)^+$.
The fifteen polynomials and their effective indices are given in
Table~\ref{tab:elliptic-values}.
For every prime $p\ge5$ with $p\ne N$, the reduction
$\overline{\mathfrak h}_N$ is squarefree, and its roots in $\Fb$ are exactly the
values with effective index $e>1$. Consequently,
\[
\gcd(\mathcal S_{N,p},\overline{\mathfrak h}_N)=\prod_{e_i>1}(T-t_i).
\]
\end{lemma}
\begin{proof}
The characteristic-zero polynomials are those of
\citet[\S4.2, Table~1]{JST19}, with $T_N=q^{-1}+O(q)$.
Their elliptic orders give the effective indices of
Lemma~\ref{lem:inertia}. For $N\in\mathcal N_0$, the Fricke fixed values
are given by $u_N^2=b_N$. At the remaining levels, the smooth models
$v^2=f_N(u)$ of Section~\ref{sec:comp} have involution
$(u,v)\mapsto(u,-v)$ and coordinate $T_N=u-s_N$;
the branch values are the roots of $f_N(T+s_N)$, all of index two.
The additional factors are $T+2$ of index two at $N=17,29$,
$T$ of index two at $N=41$, $T+3$ of index three at $N=19$,
and $T$ of index three at $N=31$. These are obtained from the
elliptic fibres over $j=1728$ and $j=0$.

We justify specialization of the support. An effective stabilizer arises
from an automorphism of $(E,C)$ or a Fricke fixed point. For $p\ge5$,
nontrivial reduced automorphisms occur only at $j=0,1728$; their
prime-to-$p$ automorphisms lift with the corresponding CM elliptic curve,
and an invariant prime-to-$p$ cyclic subgroup lifts uniquely over a
strict henselian base. Conversely, these automorphisms persist on reduction.
The fixed-point scheme of $w_N$ on $X_0(N)_{\Lambda_N}$ is finite
\'etale over $\Lambda_N$. To see this, let $x$ be a fixed geometric point
of a fibre. After passing to the strict henselization of the base and
completing, the local ring at $x$ is $\mathcal O[[z]]$, where $\mathcal O$
is a complete local ring in which $2$ is invertible. The involution acts
on the cotangent line of the fibre at $x$ by a sign. The sign is $-1$: a
tame automorphism of finite order of a smooth connected curve that fixes
a point and acts trivially on its cotangent line is the identity, whereas
$w_N$ exchanges the two cusps. Hence $z'=(z-w_N(z))/2$ is again a
parameter, $\mathcal O[[z]]=\mathcal O[[z']]$, and $w_N(z')=-z'$. The
ideal generated by all $w_N(f)-f$ is then $(2z')=(z')$, so the
fixed-point scheme is $\operatorname{Spec}\mathcal O$ near $x$, a section.
Being also proper and quasi-finite, the fixed-point scheme is finite
\'etale. Fricke fixed points therefore neither appear nor merge under
reduction at a prime not dividing $2N$.
Thus the listed elliptic values specialize to precisely the values with
$e>1$.

Exact discriminant calculation for every displayed $\mathfrak h_N$ gives
\[
\operatorname{Disc}(\mathfrak h_N)=\pm2^{\lambda_N}3^{\mu_N}N^{\nu_N}
\qquad(\lambda_N,\mu_N,\nu_N\ge0).
\]
The exact values and model calculations are recorded in the supplementary
material (Section~\ref{sec:comp}).
Hence distinct elliptic values cannot collide at any admissible prime,
so $\overline{\mathfrak h}_N$ is squarefree. Intersecting its root set with the
supersingular locus proves the gcd formula.
\end{proof}

\begin{table}[p]\centering\small
\setlength{\tabcolsep}{4pt}
\begin{tabular}{cll}\toprule
$N$ & $\mathfrak h_N(T)$ & Effective indices\\\midrule
2 & $(T+104)(T-152)$ & $4,\ 2$\\[3pt]
3 & $(T+42)(T-66)$ & $6,\ 2$\\[3pt]
5 & $(T+16)(T^{2}-12T-464)$ & $2,\ 2$\\[3pt]
7 & $(T+10)(T+9)(T-18)$ & $2,\ 3,\ 2$\\[3pt]
11 & $(T+6)(T^{3}-2T^{2}-76T-212)$ & $2,\ 2$\\[3pt]
13 & $(T+4)(T+3)(T^{2}-4T-48)$ & $2,\ 3,\ 2$\\[3pt]
17 & $(T+2)(T^{4}+2T^{3}-39T^{2}-176T-212)$ & $2,\ 2$\\[3pt]
19 & $(T+4)(T+3)(T^{3}-4T^{2}-16T-12)$ & $2,\ 3,\ 2$\\[3pt]
23 & $\begin{aligned}[t]&(T^{3}-2T^{2}-17T-25)\\[-1pt]&{}\cdot(T^{3}+6T^{2}+11T+7)\end{aligned}$ & $2,\ 2$\\[3pt]
29 & $(T+2)(T^{6}+2T^{5}-17T^{4}-66T^{3}-83T^{2}-32T-4)$ & $2,\ 2$\\[3pt]
31 & $T(T^{3}-17T-27)(T^{3}+4T^{2}+3T+1)$ & $3,\ 2,\ 2$\\[3pt]
41 & $\begin{aligned}[t]T\bigl(&T^8+4T^7-8T^6-66T^5-120T^4\\[-1pt]&{}-56T^3+53T^2+36T-16\bigr)\end{aligned}$ & $2,\ 2$\\[3pt]
47 & $\begin{aligned}[t]&(T^{5}-5T^{3}-20T^{2}-24T-19)\\[-1pt]&{}\cdot(T^{5}+4T^{4}+7T^{3}+8T^{2}+4T+1)\end{aligned}$ & $2,\ 2$\\[3pt]
59 & $\begin{aligned}[t]&(T^{3}+2T^{2}+1)\\[-1pt]&{}\cdot(T^{9}+2T^{8}-4T^{7}-21T^{6}-44T^{5}-60T^{4}-61T^{3}-46T^{2}-24T-11)\end{aligned}$ & $2,\ 2$\\[3pt]
71 & $\begin{aligned}[t]&(T^{7}-7T^{5}-11T^{4}+5T^{3}+18T^{2}+4T-11)\\[-1pt]&{}\cdot(T^{7}+4T^{6}+5T^{5}+T^{4}-3T^{3}-2T^{2}+1)\end{aligned}$ & $2,\ 2$\\[3pt]
\bottomrule\end{tabular}
\caption{Elliptic-value polynomials at all fifteen prime levels, in the
normalization $T_N=q^{-1}+O(q)$, from \citet[\S4.2, Table~1]{JST19}.
The indices follow the displayed factors; each index applies to every root
of its factor. At admissible primes, the supersingular roots of
$\overline{\mathfrak h}_N$ are exactly the thick values removed in
\eqref{eq:elliptic-gcd}.}
\label{tab:elliptic-values}
\end{table}


\begin{lemma}\label{lem:AW}
As mod-$p$ forms on $\Gamma_0(N)$: $A|_{p-1}W_N=\Lp{N}A$, $\Delta_N|_{24}W_N=\Delta_N$ and
$\theta T|_2W_N=\theta T$.
\end{lemma}
\begin{proof}
For a level-one form $f$ of weight $w$, $f|_wW_N=N^{w/2}f(N\tau)$. The form $E_{p-1}(N\tau)$ has
$q$-expansion $\equiv1\equiv E_{p-1}(\tau)$, so it equals $A$ modulo $p$ by Fact~\ref{fact:modp}(i);
and $N^{(p-1)/2}\equiv\Lp N$. The other two identities hold in characteristic $0$:
$\Delta(-1/(N\tau))\Delta(-1/\tau)=N^{12}\tau^{24}\Delta(N\tau)\Delta(\tau)$, and $T$ is
$w_N$-invariant.
\end{proof}

\begin{proposition}[weighted supersingular polynomial]\label{prop:S}
Let $S(T)=\prod_{j=1}^s(T-t_j)^{(p+1)/e_j}$. Then $S\in\Fp[T]$ and
\[
S(\overline T)=\Delta_N^{-k}=q^{-(N+1)k}\prod_{n\ge1}(1-q^n)^{-24k}(1-q^{Nn})^{-24k}\quad\text{in }\Fp((q)).
\]
Equivalently, $S=A^{p+1}/\Delta_N^{k}$ as functions on $X_0(N)^+_{\Fp}$. In particular
$\sum_j e_j^{-1}=(N+1)(p-1)/24$.
\end{proposition}
\begin{proof}
The function $\Sigma:=A^{p+1}/\Delta_N^k$ is a ratio of mod-$p$ forms of weight $p^2-1$. By Lemma~\ref{lem:AW} it is
$w_N$-invariant, since $\Lp N^{p+1}=1$. $\Delta_N$ does not vanish on $Y_0(N)$, so $\Sigma$ is regular off the cusp of
$X_0(N)^+$. Its $q$-expansion there is $\Delta_N^{-k}=q^{-(N+1)k}(1+O(q))$, because $A$ has
$q$-expansion $1$. Hence $\Sigma$ is a monic polynomial in $T$ of degree $(N+1)k$. By
Fact~\ref{fact:modp}(ii) and Lemma~\ref{lem:inertia}(ii), $\Sigma$ vanishes exactly at the supersingular
$t_j$, to order $(p+1)/e_j$ in $T-t_j$. Comparing degrees gives the mass identity.
\end{proof}

\begin{corollary}\label{cor:degree}
The degree of the reduced denominator is
\[
\deg D_{\rm thin}=\frac{(N+1)(p-1)}{24}-\sum_{e_i>1}\frac1{e_i}.
\]
In particular $R_{N,p}=0$ exactly when no thin supersingular value exists.
\end{corollary}
\begin{proof}
Separate the thin and non-thin terms in the mass identity of
Proposition~\ref{prop:S}; use Theorem~\ref{thm:main}(b) for the last assertion.
\end{proof}

\section{The Fricke residue formula}\label{sec:kz}
We now apply the arithmetic identities to the Fricke setting of
Section~\ref{sec:intro}. The derivative argument gives the pole support
and local residues; Theorem~\ref{thm:input} determines the signed
weighted product.

\subsection{The derivative and pole support}

The following is the level-$N$ version of \citet[\S10, (35)--(38)]{KZ98}, who credit the method to Koike.

\begin{lemma}\label{lem:dpsi}
In $\Fp((q))$, with $d/dT=(\theta T)^{-1}\theta$,
\[
\frac{d\overline\psi}{dT}=T^{p-1}-(\theta T)^{p-1},\qquad (\theta T)^{p-1}=g(\overline T),\quad
g:=\frac{(\theta T)^{p-1}}{A^2}\in\Fp(T).
\]
\end{lemma}
\begin{proof}
Over $\Lambda_N$, $\theta\psi=T^{p-1}\theta T-(\theta T)(q^p)$, and $(\theta T)(q^p)\equiv(\theta T)^p$.
Divide by $\theta T=-q^{-1}+\cdots$, which is invertible. The function $g$ is a ratio of mod-$p$
forms of weight $2(p-1)$, it is $w_N$-invariant by Lemma~\ref{lem:AW}, and it has $q$-expansion
$(\theta T)^{p-1}$ because $A\equiv1$.
\end{proof}

\begin{lemma}[local structure of $g$]\label{lem:g}
Let $t\in\Fq$ be a non-cuspidal point with effective inertia $e_t$, and put $\varepsilon_t=1$ if
$t$ is supersingular and $0$ otherwise. Then
\[
\mathrm{ord}_{T=t}\,g=\frac{(p-1)(e_t-1)-2\varepsilon_t}{e_t}.
\]
For thin $t$ and $x$ over $t$, let $a=A/(\theta T)^{(p-1)/2}$, a function on $X_0(N)$ near $x$.
Then $a$ has a simple zero at $x$,
\[
a'(t):=\frac{da}{dT}(x)=\frac{(B/12)(x)}{(\theta T)(x)^{(p+1)/2}}\ne0,
\qquad g=a'(t)^{-2}(T-t)^{-2}\bigl(1+O(T-t)\bigr).
\]
\end{lemma}
\begin{proof}
By Fact~\ref{fact:modp}(iv) and Lemma~\ref{lem:inertia}(ii), $\theta T$ vanishes to order $e_t-1$ in a
stack parameter $z$, and $A$ to order $\varepsilon_t$; also $T-t\sim z^{e_t}$. This gives the order formula.
For thin $t$, the map $X_0(N)\to X_0(N)^+$ is \'etale at $x$ and $\theta T(x)\ne0$. In
characteristic $0$, for forms $f$ and $M$ of equal weight $w$,
$\theta(f/M)=(\partial f\cdot M-f\cdot\partial M)/M^2$, where $\partial=\theta-\frac{w}{12}E_2$ is the
Serre derivative. Apply this with $f=E_{p-1}$ and $M=(\theta T)^{(p-1)/2}$ and reduce modulo $p$; at
$x$ we have $A(x)=0$. The form $\partial E_{p-1}$ has weight $p+1$ and $q$-expansion
$\equiv\theta(1)+\frac1{12}E_2\equiv\frac1{12}E_{p+1}$, using the Kummer congruence
$2(p+1)/B_{p+1}\equiv24$ and $\sigma_p\equiv\sigma_1$. So $\partial A=B/12$ as mod-$p$ forms. Finally
$da/dT=\theta a/\theta T$. The value $B(x)\ne0$ because $B/12$ is the derivative of $A$ at its simple zero.
\end{proof}

\begin{proof}[Proof of Theorem~\ref{thm:main}(b)--(c)]
Let $t\in\Fq$. By Proposition~\ref{prop:res}, near $t$ we have
$\overline\psi=-\overline H(t^p,t)(T-t)^{-1}+O(1)$, because $T^{p^2}-T=-(T-t)(1+O(T-t))$. So the
coefficient of $(T-t)^{-2}$ in $d\overline\psi/dT$ is $\overline H(t^p,t)$. By Lemma~\ref{lem:dpsi}
it is also $-[g]_{-2}$. Hence
\[
\Res_tR_{N,p}=-\overline H(t^p,t)=[g]_{-2}.
\]
By Lemma~\ref{lem:g}, $\mathrm{ord}_tg\ge0$ unless $t$ is thin. For $p\ge5$ and $e_t\ge2$ we have
$(p-1)(e_t-1)-2\ge0$. At thin $t$, $[g]_{-2}=a'(t)^{-2}=144(\theta T)(x)^{p+1}/B(x)^2\ne0$. Every pole is
simple by Proposition~\ref{prop:res}, so $R_{N,p}=P/D_{\rm thin}$ with $P(t_i)\ne0$ at every thin
$t_i$. Since $R_{N,p}(\infty)=0$, $\deg P<\deg D_{\rm thin}$.
\end{proof}

\subsection{The weighted signed product}

\begin{proposition}[weighted residue product]\label{prop:product}
For thin $t_i$ and $x=(E,C)$ over $t_i$,
\[
r_iW_i=\frac{B^{p-1}}{\Delta_N^{\,k}}(x)=\beta(E)\,D_N(x)^k.
\]
\end{proposition}
\begin{proof}
Near $x$, Proposition~\ref{prop:S} gives $S(T)=A^{p+1}/\Delta_N^k=a^{p+1}\,\xi^{k}$ with
$\xi:=(\theta T)^{12}/\Delta_N$, a function that is regular and nonzero at $x$. Hence
$W_i=[S/(T-t_i)^{p+1}]_{T=t_i}=a'(t_i)^{p+1}\xi(x)^k$. With $r_i=a'(t_i)^{-2}$ this gives
\[
r_iW_i=a'(t_i)^{p-1}\xi(x)^k=\Bigl(\frac{B}{12}\Bigr)^{p-1}(\theta T)^{-\frac{p^2-1}2}\cdot
\frac{(\theta T)^{\frac{p^2-1}{2}}}{\Delta_N^{k}}=\frac{B^{p-1}}{\Delta_N^k}(x).
\]
Since $\Delta_N=\Delta(\tau)\Delta(N\tau)=\Delta^2/D_N$, we have
$\Delta_N^k=\Delta^{2k}D_N^{-k}$, which gives the second equality.
\end{proof}

\begin{proof}[Proof of Theorem~\ref{thm:main}(e)]
Proposition~\ref{prop:product} and Theorem~\ref{thm:input} give directly
\[
r_iW_i=\beta(E)\,D_N(x)^k=-\Lp{-1}\Lp N=-\Lp{-N}.
\]
By Corollary~\ref{cor:rFp}, $r_i\in\Fp^\times$, and hence
$W_i=-\Lp{-N}r_i^{-1}\in\Fp^\times$.
\end{proof}

\begin{proof}[Proof of Theorem~\ref{thm:main}(f)]
Proposition~\ref{prop:res} gives
$\overline H(t^p,t)=-\Res_tR_{N,p}$. By part (b) this is zero
unless $t=t_i$ is thin, in which case it equals
$-r_i=\Lp{-N}W_i^{-1}$ by part (e).
\end{proof}

\begin{remark}[$\Fp$-rationality at non-rational points]
The two rationality assertions in Theorem~\ref{thm:main}(e) have different
explanations. The assertion $r_i\in\Fp$ follows from symmetry of the
Hecke polynomial and Frobenius-equivariance of residues
(Corollary~\ref{cor:rFp}). In contrast, writing
\[
W_i=\operatorname{Nm}_{\Fq/\Fp}(D_{\rm thin}'(t_i))
\prod_{e_j>1}(t_i-t_j)^{(p+1)/e_j},
\]
shows that $W_i\in\Fp$ requires the product of the thick-point factors
to lie in $\Fp$. This is not visible factor by factor; it follows here
from the signed product identity together with Corollary~\ref{cor:rFp}:
$W_i=-\Lp{-N}/r_i\in\Fp^\times$.
\end{remark}


\section{Geometric interpretation of the weights}\label{sec:geometry}
The proof of Theorem~\ref{thm:main} is algebraic. We now interpret its weights in
the integral modular curve, retaining the geometric explanation of pole cancellation.
Nothing in Sections~\ref{sec:imports}--\ref{sec:kz} depends on this section, and the
argument given for Proposition~\ref{prop:geometry} is an outline that rests on the
cited results of Deligne--Rapoport and Liu.
Let $\mathbb W$ be the ring of Witt vectors of $\Fb$ and let
\[
\mathcal B=X_0(Np)^{\mathrm{DR}}_{\mathbb W},\qquad
\mathcal A=\mathcal B/\langle w_N\rangle.
\]
Here $\mathcal B$ denotes the coarse Deligne--Rapoport model and $w_N$ is the
partial Atkin--Lehner involution. A split node with completed local ring
$\mathbb W[[s,t]]/(st-p^ev)$, $v$ a unit, has \emph{thickness} $e$.

\begin{proposition}\label{prop:geometry}
The special fibre of $\mathcal A$ has two components, each isomorphic to
$X_0(N)^+_{\Fb}$. They meet at the supersingular Fricke orbits. At the orbit of
$x=(E,C_N)$ the node has thickness
\[
e_x=|\Aut(E,C_N)/\{\pm1\}|\,
|\Stab_{\langle w_N\rangle}(x)|.
\]
Consequently $R_{N,p}$ has a pole at a supersingular value precisely when the
corresponding node has thickness one.
\end{proposition}
\begin{proof}[Outline of proof]
Deligne--Rapoport's theorem applies with prime-to-$p$ level $N$ and the stabilizer
of a line in $(\Z/N)^2$ as the auxiliary group
\citep[V, 1.14--1.20; VI, Theorem~6.9, especially (6.9.1)]{DR73}.
It gives two copies of $X_0(N)_{\Fb}$ with Frobenius gluing, and the coarse
completed local ring at $x$ is
\[
\mathbb W[[s,t]]/(st-p^{m_x}),\qquad m_x=|\Aut(E,C_N)/\{\pm1\}|.
\]
The forgetful map $\mathcal B\to X_0(N)_{\mathbb W}$ is finite: it is proper and has finite
fibres on the generic fibre and on the two special-fibre components.
Normality identifies $\mathcal B$ with the normalization in its generic function
field. The generic partial Fricke map covers $w_N$ on $X_0(N)_{\mathbb W}$, hence extends
uniquely by normalization. On the elliptic locus it is
\[
(E,C_N,C_p)\longmapsto(E/C_N,E[N]/C_N,q_N(C_p)).
\]
Since $q_N$ is prime to $p$, it preserves the Frobenius and Verschiebung components
and both branches at a node.

To apply Liu's quotient result, note that $\mathcal B$ is a normal,
proper flat fibred surface over the DVR $\mathbb W$, with smooth generic fibre
and nodal special fibre. At a fixed geometric node its stabilizer has
order two, prime to $p$, and acts trivially on its residue field $\Fb$.
As just shown it preserves both branches. It acts nontrivially on either
coarse branch. Indeed its action on the smooth curve $X_0(N)$ is a nontrivial
tame involution, whose derivative at a fixed point is $-1$.
Taking the branch-preserving quotient therefore multiplies the thickness by
this order, by \citet[Proposition~10.3.48(b)--(c), pp.~526--528]{Liu02}.
At a non-fixed node the inertia is trivial and the thickness is unchanged.
Invariants commute with reduction because $2$ is invertible in $\mathbb W$.
The quotient components are consequently $X_0(N)^+_{\Fb}$, and no node is
smoothed, since no stabilizer exchanges branches. This proves the thickness formula.
The last assertion is Theorem~\ref{thm:main}(b).
\end{proof}

\begin{remark}[level two]
For $N=2$, the Fricke fixed points are $u=\pm64$, with $T=-104,152$.
At $u=-64$ the reduced automorphism group has order two and the Fricke stabilizer
has order two, giving thickness four. At $u=64$ their orders are one and two,
giving thickness two. All other supersingular values have thickness one.
The two exceptional values are supersingular precisely when
$\left(\frac{-1}{p}\right)=-1$ and $\left(\frac{-2}{p}\right)=-1$, respectively.
Thus a thick value is absent from the pole divisor but remains in the
weighted residue product. The same distinction applies at all fifteen levels.
\end{remark}

\section{Finite certificates and computational evidence}\label{sec:comp}
For Theorem~\ref{thm:main}, exact-arithmetic computations cover all fifteen
levels in $\mathcal N$ for $5\le p\le47$, excluding $p=N$. They reconstruct
$R_{N,p}$ independently of the signed residue formula and then verify
that formula. They are checks on the theorems, not inputs to their proofs.
The following bound specifies the strength of the certificates.

\begin{proposition}[finite certificate]\label{prop:certificate}
Let $L=p^2$, let $D\in\Fp[T]$ be a monic divisor of $T^L-T$, of degree $d$, and
let $P\in\Fp[T]$ have degree less than $d$. If
\[
D(\overline T)\,\overline{T|U_p}-P(\overline T)=O(q^{L-d+1}),
\]
then $R_{N,p}=P/D$. For $D=1$, take $P=0$.
\end{proposition}
\begin{proof}
Proposition~\ref{prop:res} implies that
$R_{N,p}-P/D=C/(T^L-T)$ with $\deg C<L$.
If $C\ne0$, its $q$-order is $L-\deg C$, since $T=q^{-1}+O(q)$.
Multiplication by $D(T)$ gives order $L-\deg C-d\le L-d$, contradicting the
hypothesis. Thus $C=0$.
\end{proof}

For each pair, the supersingular values and effective indices are determined
from the modular model, giving $D_{\rm thin}$. A separate $q$-expansion
calculation recovers the numerator without using the residue signs and
checks Proposition~\ref{prop:certificate}, using source precision at least
$p^3+2$. The reduced fraction, mass identity, weighted supersingular
polynomial and signed residue products are also checked.

At $N\in\mathcal N_0$, the inputs are the eta coordinate and its rational
$j$-map. At the other ten levels, we use the standard smooth elliptic or
hyperelliptic models $v^2=f_N(u)$ with $w_N(u,v)=(u,-v)$ and relations
\[
j+j\circ w_N=\alpha_N(u),\qquad
j(j\circ w_N)=Q_N(u),
\]
\[
j=\frac{\alpha_N(u)+\beta_N(u)v}{2},
\quad Q_N=\frac{\alpha_N^2-f_N\beta_N^2}{4}.
\]
These are the model data used by SageMath's prime-isogeny routines
\citep{SageModels}. The coordinate is $T_N=u-s_N$, where
$s_N=-[u^{N-1}]\alpha_N/N$. Formal Newton inversion of
$\alpha_N(u(q))=j(q)+j(q^N)$ gives its expansion, with the product relation
checked separately. Supersingular fibres satisfy
$j_0^2-j_0\alpha_N(u)+Q_N(u)=0$; the roots of $f_N$ detect Fricke
ramification, and fibre multiplicities over $j_0=0,1728$ give the reduced
automorphism orders. Thus the indices are obtained independently of both
the signed residue formula and the weighted Hasse polynomial.

All 184 admissible pairs pass the finite certificate. They contain
4458 thin values, of which 3448 lie outside $\Fp$
(Table~\ref{tab:all-levels}).
At every thin value the weighted residue product is
$-\left(\frac{-N}{p}\right)$; pairs with no thin values have an empty
residue-product set.

\begin{table}[htbp]\centering\small
\begin{tabular}{rrrr}\toprule
$N$ & certified pairs & thin values & outside $\Fp$\\\midrule
2&13&33&8\\
3&13&47&16\\
5&12&67&32\\
7&12&92&54\\
11&12&136&86\\
13&12&163&116\\
17&12&203&140\\
19&12&231&170\\
23&12&267&202\\
29&12&334&260\\
31&12&350&272\\
41&12&445&352\\
47&12&483&382\\
59&13&733&614\\
71&13&874&744\\
\midrule
Total&184&4458&3448\\\bottomrule
\end{tabular}
\caption{Certified reconstructions at all fifteen levels for
$5\le p\le47$, $p\ne N$. Values are distinct geometric $T_N$-values
within each pair $(N,p)$; counts are summed over the tested primes.}
\label{tab:all-levels}
\end{table}

Further tests were run at the five eta-coordinate levels $N\in\mathcal N_0$.
All 222 admissible pairs through $p=211$ are certified
(6246 thin values, 4696 outside $\Fp$); recovery of the denominator without
a candidate was checked through $p=37$, and Hecke polynomials modulo $p^2$
through $p=31$. A further 960 pairs for $223\le p\le1499$ were checked using the
double-pole formula of Lemma~\ref{lem:g}; these are not independent
reconstructions from $T|U_p$.

\begin{remark}[size of the input data]\label{rem:datasize}
Corollary~\ref{cor:reconstruction} reconstructs $R_{N,p}$ from the
supersingular values and their effective indices. By the mass identity of
Proposition~\ref{prop:S} and Lemma~\ref{lem:elliptic-specialization}, there
are at most $(N+1)(p-1)/24+\deg\mathfrak h_N$ such values. By contrast, the
certificate of Proposition~\ref{prop:certificate} uses the coefficients of
$T|U_p$ up to $q^{p^2-d}$, where $d=\deg D_{\rm thin}$, that is, those of $T$ up to $q^{p(p^2-d)}$.
This compares the size of the input only. Obtaining the supersingular
values has its own cost, and no claim about running times is made here;
timings for particular implementations are recorded in the supplementary
material.
\end{remark}

\paragraph{Data and software.}
The supplementary material for this paper consists of five packages.
Their folder names are those of the archived release; they do not reflect
the current status or numbering of the results.
\begin{itemize}[leftmargin=*,itemsep=1pt]
\item \path{Signed_Weighted_Residue_v1_0}: code, per-pair JSON results
and reproduction instructions for the residue formula at
$N\in\mathcal N_0$ (\texttt{python-flint} 0.9.0 and \texttt{mpmath}).
\item \path{Fricke_Fifteen_Levels_Extension_v1_0}: the verifier for the
ten levels outside $\mathcal N_0$, a regression run at the five levels in
$\mathcal N_0$, the model coefficients with their provenance, and
per-pair certificates (\texttt{python-flint} 0.9.0).
\item \path{Fricke_Section8_Unified_Certificates_v1_0}: the records and
reproduction script for Table~\ref{tab:all-levels}.
\item \path{Fricke_Elliptic_Values_Fifteen_Levels_v1_0}: the
discriminants and model calculations for
Lemma~\ref{lem:elliptic-specialization}.
\item \path{Corollary_1_2_Reconstruction_Benchmark_v1_0}: timings for the
reconstruction of Corollary~\ref{cor:reconstruction} (numbered
Corollary~1.2 in an earlier draft) at $N\in\mathcal N_0$, with the timing
code, exact inputs and outputs, raw timing samples and verification
records.
\end{itemize}
\if\relax\detokenize\expandafter{\DataRepositoryURL}\relax
The code, per-pair data and reproduction instructions are supplied as
supplementary files with the manuscript.
\else
The code, data and reproduction instructions are publicly available at
\expandafter\url\expandafter{\DataRepositoryURL}.
\fi

\section{Further directions}
The discriminant-ratio identity of Theorem~\ref{thm:input}(ii) holds without
a genus-zero hypothesis, including for arbitrary prime-to-$p$ isogenies
\citep{JordanFactorial}. The residue formula additionally uses an integral genus-zero quotient coordinate,
the Hasse divisor and the Kodaira--Spencer derivative.
Lemma~\ref{lem:coordinate} supplies this coordinate uniformly for all
fifteen prime Fricke levels; a coordinate on $X_0(N)$ is unnecessary.
At composite levels, the cusp structure and the chosen Atkin--Lehner
quotient require separate analysis. For a quotient of positive genus,
there is no global Hauptmodul, and the residue reconstruction must be
reformulated geometrically. These extensions of the residue theorem
remain research questions; neither is established here.

\section*{Code and data availability}
All code, scripts and logs are available in the GitHub repository
\begin{center}\url{https://github.com/bbrhuft/fricke-signed-residues}\end{center}
The computational materials described here are those of release v2.0 of this repository.
Earlier versions of this work, in which this paper and its companion formed a single
manuscript, are archived on Zenodo at
\href{https://doi.org/10.5281/zenodo.24024359}{doi:10.5281/zenodo.24024359} (v1.11) and
\href{https://doi.org/10.5281/zenodo.24084789}{doi:10.5281/zenodo.24084789} (v1.13).

\begin{thebibliography}{99}\small
\bibitem[Baker(1998)]{Baker98} Baker, A. (1998). A supersingular congruence for modular forms.
\emph{Acta Arithmetica}, 86(1), 91--100. \url{https://doi.org/10.4064/aa-86-1-91-100}.


\bibitem[Betina and Lecouturier(2018)]{BL18} Betina, A., \& Lecouturier, E. (2018). Congruence formulas for
Legendre modular polynomials. \emph{Journal of Number Theory}, 188, 71--87.


\bibitem[Bringmann et~al.(2026)]{BGKMR26} Bringmann, K., Guerzhoy, P.,
Kane, B., Mertens, M., \& Rolen, L. (2026). On $U_p$-congruences for
meromorphic modular forms via supersingularity. \emph{arXiv:2606.14020v1}.
\url{https://arxiv.org/abs/2606.14020v1}.


\bibitem[Chen et~al.(2019)]{CMT19} Chen, R.~C., Marks, S., \& Tyler, M. (2019).
$p$-Adic properties of Hauptmoduln with applications to moonshine.
\emph{SIGMA}, 15, 033, 35 pp. \url{https://doi.org/10.3842/SIGMA.2019.033}.


\bibitem[Conway and Norton(1979)]{CN79} Conway, J.~H., \& Norton, S.~P. (1979). Monstrous moonshine.
\emph{Bulletin of the London Mathematical Society}, 11, 308--339.


\bibitem[Deligne and Rapoport(1973)]{DR73} Deligne, P., \& Rapoport, M. (1973). Les sch\'emas de modules de
courbes elliptiques. In \emph{Modular Functions of One Variable II} (LNM 349, pp.~143--316). Springer.


\bibitem[Deuring(1941)]{Deuring41} Deuring, M. (1941). Die Typen der Multiplikatorenringe elliptischer Funktionenk\"orper. \emph{Abhandlungen aus dem Mathematischen Seminar der Hansischen Universit\"at}, 14, 197--272.


\bibitem[Duncan and Swisher(2026)]{DS26} Duncan, J.~F.~R., \& Swisher, H. (2026).
Modular functions and the monstrous exponents. \emph{arXiv:2602.09135v1}, 9~February~2026. \url{https://arxiv.org/abs/2602.09135v1}.


\bibitem[Gross(1990)]{Gross90} Gross, B.~H. (1990). A tameness criterion for Galois representations associated to modular forms (mod $p$). \emph{Duke Mathematical Journal}, 61(2), 445--517.


\bibitem[Jordan(2026)]{JordanFactorial} Jordan, D. (2026). A factorial supersingular congruence and
signed discriminant ratios. Preprint (companion paper). Manuscript, code and data: \url{https://github.com/bbrhuft/factorial-supersingular-congruence}.

\bibitem[Jorgenson et~al.(2019)]{JST19} Jorgenson, J., Smajlovi\'c, L., \& Then, H. (2019).
The Hauptmodul at elliptic points of certain arithmetic groups.
\emph{Journal of Number Theory}, 204, 661--682.
\url{https://doi.org/10.1016/j.jnt.2019.03.021}.
The elliptic-value tables used here are also available in
\emph{arXiv:1602.07426v2}, \S4.2, Table~1 and Appendix~C.
\url{https://arxiv.org/abs/1602.07426v2}.


\bibitem[Kaneko and Zagier(1998)]{KZ98} Kaneko, M., \& Zagier, D. (1998). Supersingular $j$-invariants,
hypergeometric series, and Atkin's orthogonal polynomials. In \emph{Computational Perspectives on Number
Theory} (AMS/IP Studies in Advanced Mathematics 7, pp.~97--126).


\bibitem[Katz and Mazur(1985)]{KM85} Katz, N.~M., \& Mazur, B. (1985). \emph{Arithmetic Moduli of Elliptic
Curves} (Annals of Mathematics Studies 108). Princeton University Press.


\bibitem[Katz(1973)]{Katz73} Katz, N.~M. (1973). $p$-adic properties of modular schemes and modular forms.
In \emph{Modular Functions of One Variable III} (LNM 350, pp.~69--190). Springer.


\bibitem[Koike(1973)]{Koike73} Koike, M. (1973). Congruences between modular forms
and functions and applications to the conjecture of Atkin.
\emph{Journal of the Faculty of Science, University of Tokyo, Section IA, Mathematics},
20, 129--169.


\bibitem[Liu(2002)]{Liu02} Liu, Q. (2002). \emph{Algebraic Geometry and Arithmetic Curves}. Oxford Graduate Texts in Mathematics 6. Oxford University Press.


\bibitem[Morton(2023)]{Morton23} Morton, P. (2023). Supersingular conjectures for the
Fricke group. \emph{International Journal of Number Theory}, 19(4), 803--828.
\url{https://doi.org/10.1142/S1793042123500409}; preprint
\url{https://arxiv.org/abs/2105.14532}.


\bibitem[Nakaya(2019)]{Nakaya19} Nakaya, T. (2019). The number of linear factors of
supersingular polynomials and sporadic simple groups. \emph{Journal of Number Theory},
204, 471--496. \url{https://doi.org/10.1016/j.jnt.2019.04.015}.


\bibitem[Ogg(1975)]{Ogg75} Ogg, A.~P. (1975).
Automorphismes de courbes modulaires.
\emph{S\'eminaire Delange--Pisot--Poitou, Th\'eorie des nombres},
16 (1974--1975), no.~1, expos\'e~7, 1--8.
\url{https://www.numdam.org/item/SDPP_1974-1975__16_1_A4_0/}.


\bibitem[Ono(2026)]{Ono25} Ono, K. (2026). Modular forms for chromatic homotopy: supersingular congruences. \emph{arXiv:2509.16175v2}, 3 July 2026. \url{https://arxiv.org/abs/2509.16175v2}.


\bibitem[SageMath Developers(2026)]{SageModels} The SageMath Developers (2026).
\emph{Isogenies of small prime degree}, including
\texttt{\_hyperelliptic\_isogeny\_data}, in the SageMath source code.
\url{https://github.com/sagemath/sage/blob/develop/src/sage/schemes/elliptic_curves/isogeny_small_degree.py}.
The exact coefficient data and source checksum used here are preserved
in the supplementary extension package.


\bibitem[Sakai(2011)]{Sakai11} Sakai, Y. (2011). The Atkin orthogonal polynomials for
the low-level Fricke groups and their application.
\emph{International Journal of Number Theory}, 7(6), 1637--1661.
\url{https://doi.org/10.1142/S1793042111004460}.


\bibitem[Sakai(2014)]{Sakai14} Sakai, Y. (2014). The Atkin orthogonal polynomials for
the Fricke groups of levels $5$ and $7$.
\emph{International Journal of Number Theory}, 10(8), 2243--2255.
\url{https://doi.org/10.1142/S1793042114500766}.


\bibitem[de~Shalit(1994)]{dS94} de~Shalit, E. (1994). Kronecker's polynomial, supersingular elliptic curves,
and $p$-adic periods of modular curves. \emph{Contemporary Mathematics}, 165, 135--148.


\bibitem[Waterhouse(1969)]{Waterhouse69} Waterhouse, W.~C. (1969). Abelian varieties over finite fields. \emph{Annales Scientifiques de l'\'Ecole Normale Sup\'erieure (4)}, 2, 521--560.
\end{thebibliography}
	\bigskip
\noindent{\scshape Maynooth University Social Sciences Institute (MUSSI), Maynooth University, Maynooth, County Kildare, Ireland}\\
\emph{Email address}: \texttt{David.Jordan@mu.ie}
\end{document}

